% Source: 24_Chapter_16_Bending_the_Necks_of_Victors_Mission_to.txt
\bookchapter[Bending the Necks of Victors]{CHAPTER 16}{BENDING THE NECKS OF VICTORS}{Mission to the Barbarians}

\leadin{T}{HE BRISTLING FORESTS OF} northern Europe were inhabited by barbarian tribes who
sacrificed animals and worshiped nature spirits among the trees and beside the streams. Any
missionary who ventured among them, with any hope of conversions, had to demonstrate the superior
power of the Christian spirits.

The most famous incident tells of one eighth-century missionary named Boniface who marched into a
shrine in Germany, the sacred forest of Thor, the god of thunder. The cult object was a massive oak.
Boniface, so the story goes, took an axe to it. Just as he leveled the first stroke, a mighty breath
of wind from God toppled the tree. The pagans marveled and were converted. Boniface used the wood to
build a chapel to St. Peter.

That is the stuff that made Europe Christian. The missionary monks overthrew barbarian magic by
calling down superior powers. God blew, and the tree fell. A miracle here and a victorious battle
there; and the Germans were ready for baptism.

The point is, Christianity faced a new and enormous challenge when the German tribes swept across
the Rhine and Danube frontier and left the once mighty empire of the Romans in shambles. The
civilization that arose on those ruins was Christian in character because the invaders surrendered,
not their arms, but their gods.

How did the barbarians come to accept the chapel dedicated to St. Peter where the sacred oak of Thor
once stood? Why did the Germans come to worship the God of their defeated foes?

\section{THE COMING OF THE BARBARIANS}

The year 476 usually marks the end of the Christian Roman Empire in the West. That is the year the
long line of emperors inaugurated by Augustus (27 BC–AD 14) ended, and the undisguised rule by
German leaders began.

The event was about as unspectacular as events can be. The mercenary soldiers wielding the real
power in the name of the emperor were mainly German. Germans were killing Germans.

In 475 Orestes, the barbarian commander of the imperial troops, forced the Senate to elect his young
son Romulus Augustulus (Little Augustus) as emperor in the West. In the following year, however,
another Germanic commander, Odovacar, slew Orestes and, seeing no reason for continuing the sham of
an imperial line, he deposed Romulus Augustulus and proclaimed himself head of the government. That
was it. No one was particularly shocked. That had all come a generation earlier with Alaric and
Gaiseric and Attila.

Who were these new masters of Europe? The Romans had called them “barbarians” because during early
contact with the Romans they spoke no Greek and no Latin. But for the most part they were tribes
from the north, originally in or near Scandinavia: Vandals, Franks, Angles, Saxons, Goths, Lombards,
Burgundians, and others.

When they started to give the Romans problems in the third century they were at a cultural stage
midway between a pastoral and an agricultural economy. They engaged in so little commerce that
cattle, rather than money, marked the man of wealth and power among them.

According to the Roman historian Tacitus, the Germans were notoriously heavy drinkers and gamblers.
On the other hand, Tacitus praised their courage, respect for women, and freedom from many Roman
vices. A favorite amusement was listening to the tribal bards recite ancient tales of heroes and
gods. Today the tales are long lost, but we retain Thor’s day in our Thursday and another god’s
name, Wodin, in our Wednesday.

The Germans lived to fight. Every warlord had a retinue of warriors linked to their leader by a
personal oath of loyalty. According to Tacitus, any warrior leaving a battle alive after his chief
had fallen was consigned to lifelong infamy. “The chiefs fight for victory, the companions for their
chief.”

During the centuries that the Romans and Germans faced each other across the Rhine-Danube frontier,
they had countless contacts with each other, peaceful as well as hostile. Roman trade reached into
Germany, and Germans entered the empire as slaves. During the troubled third century, many
barbarians were invited to settle on vacated lands within the empire or to serve in the Roman
legions. By the end of the fourth century the Roman army and its generals in the West had become
almost completely German.

The crisis of mass invasions came with the sudden appearance of a new uncontrollable force: the
Huns. Late in the fourth century, this wave of Asiatic people crossed the Volga and soon subjugated
the easternmost Germanic tribe, the Ostragoths. Terrified at the prospect of conquest by the
advancing Huns, the Visigoths (or West Goths) petitioned the Romans to allow them to settle as
allies inside the empire. Rome granted the request, and in 376 the entire tribe crossed the Danube
into Roman territory. Shortly, however, corrupt Roman officials mistreated the Visigoths, and the
proud barbarians went on a rampage. The inept Eastern emperor Valens tried to quell them, but in 378
he lost both his army and his life in the battle of Adrianople (in today’s Turkey).

Some historians regard Adrianople as one of history’s decisive battles: it destroyed the legend of
the invincibility of the Roman legions and ushered in a century and a half of chaos. For a few years
the able emperor Theodosius I held back the Visigoths, but after his death in 395 they began to
migrate and pillage under their leader, Alaric. He invaded Italy, and in 410 his followers sacked
Rome.

To counter Alaric’s threat to Italy, the Romans had withdrawn most of their troops from the Rhine
frontier in 406 and from Britain the following year. The momentous consequence of this action was a
flood of Germanic tribes across the defenseless frontiers. The Vandals pushed their way through Gaul
to Spain and, after pressure from the Visigoths, moved on to Africa. In 455 a Vandal raiding force
sailed over from Africa, and Rome was sacked a second time.

Meanwhile the Burgundians settled in the Rhone valley, the Franks gradually spread across northern
Gaul, and the Angles, Saxons, and Jutes invaded Britain. Although each of these several tribes set
up a German-ruled kingdom within the confines of the empire, only the Franks in Gaul and the Angles
and Saxons in Britain managed to perpetuate their kingdoms longer than a few centuries.

Today we can see that the Germanic invasions were not as cataclysmic as men once thought. The
invaders seized a great deal of land, but most of this was either vacant or belonged to the
emperors; few private landowners were displaced. In most areas the Germans still represented a
minority of the population, and a gradual blending and fusing of the cultures and the blood of the
two peoples began. Thus the barbarians in time lost their Germanic customs, religion, and speech.
That is why hardly a trace of the Germanic languages remains in Italy, France, and Spain.

\section{MISSIONS TO THE NORTH}

The task of converting these northern peoples was enormous. To bring them to a nominal adherence to
Christianity was not so difficult, because they wanted to enter into the grandeur that was Rome.
Christianity was, in their eyes, the Roman religion. But to tame, refine, and educate these peoples;
to transmit to them the best of the culture of antiquity; to teach them the Christian creed; and,
above all, to instill in them even a modicum of Christian behavior—all that was another matter.

The barbarians were brought to orthodox Christianity in two ways: directly from paganism and
indirectly through Arianism, a heresy that denied the eternality of Jesus Christ as the Son of God.
Arius’s doctrine spread among them from the time in the fourth century when Arianism was strong in
the empire. This Arianism, however, was more ecclesiastical than theological, because the Germans
were not interested in the niceties of abstract theology. Their Christ was the firstborn creature.
They tended to think of him as a glorified warlord. But the main difference between the Arians and
the orthodox in the West lay in the structure of the church. The Arians had no ecclesiastical
center. They did not recognize orthodox Rome, and they had no counterpart of their own. Their
churches belonged to the clan. Even after the conversion of these Arian Germans to orthodoxy, they
were still loath to accept centralization through Rome.

The Arian influence apparently began with missionary work among the Visigoths. A half century after
the Council of Nicaea (325), a missionary named Ulfilas (or Wulfila) crossed the Danube to work
among them for forty years. Aided by other unknown missionaries, Ulfilas led them to faith in Christ
as Arius had preached him. This conversion of the Goths must have been something more than surrender
to social pressures. We know that Ulfilas translated the Bible into the Gothic language, with the
exception of the books of Kings, which he thought too bellicose to be edifying for a people already
extremely fierce and warlike. From the Goths, then, Arian Christianity spread to other German
tribes.

Some of the earliest northern peoples to come to Christ were beyond the borders of the empire and
were not themselves German. The Irish were Celtic people, and their conversion is traced to Patrick
early in the fifth century.

From his brief autobiography we learn that when the Roman legions were withdrawn for the defense of
the Continent, the Irish, then called Scots, began swooping down on the English coast, sailing up
the rivers, raiding the settlements, and carrying off plunder and slaves. Among the captives was
Patrick. So Ireland’s patron saint was not Irish! He had been reared a Christian. His father was a
deacon, but Patrick’s religion sat lightly until, as a swineherd, he prayed ardently for his
freedom. His conversion dates from this captivity. “The Lord opened to me the sense of my unbelief,”
he says, then. After six years he managed to escape and found his way to the coast where a ship
carrying a cargo of hounds was about to sail to France or Scotland. We don’t know which. Stories
about his trips to France may be sheer fancy, but we do know that he was eager to see his family,
and after many hardships reached home once again.

He would have gladly remained in England had he not had a dream one night in which the babies of
Ireland pleaded with him to come back to their country and tell them about Christ. Patrick decided
to return, but first he had to learn more about Christianity. At length he was sent out, some say
with the approval of the pope, to be a missionary to the people among whom he had once been a slave.

At this point his account ends, and from then on we have only legends. We know, however, that a
century later the structure of the church in Ireland was predominantly monastic. Presumably the
monastic community, maintaining itself on the land, fitted the agricultural communities of the Celts
better than the parish-church system so common in the Roman Empire.

We also know that Ireland became the base for the evangelization of Britain. The Irish have a habit
of leaving Ireland. And the missionary monks were no exception. As we have seen, there were
Christian churches in Britain before Patrick’s day, but a century after Patrick’s time, an Irish
monk named Columba led in the founding of a monastery on Iona, an island off the coast of Scotland.
Iona, in turn, gave a new and vital impulse to the preaching of Christ in Britain.

Surprisingly, Celtic monks from Ireland and Britain also became missionaries to the Continent. Their
itching feet carried them farther and farther from home. They established monasteries in Germany and
Switzerland and even in northern Italy. These became centers of evangelization and devout learning.
Columbanus is the great name in this mission, though he was only one among many; and St. Gall and
Bobbio were two of the most famous Celtic monasteries. These were the days before the bishops of
Rome gained general recognition as popes, so this Celtic influence on the Continent and the
independent spirit of the Irish and British monks became in time a disturbing factor to Roman
Christianity.

\section{THE GOSPEL AMONG THE FRANKS}

Roman Catholic influence began in the northern half of Gaul, among the Franks, the one tribe
destined to greatness in the shaping of Christian Europe. The founder of the nation was Clovis
(481–511), who was the first barbarian chief of any importance to convert to orthodox Christianity.
It happened that he was married to a Burgundian princess who was a Christian. Clothilda often spoke
to Clovis about the one God who created heaven and earth out of nothing and had fashioned man.
Clovis replied, “Nonsense!” But when their first son was born, he allowed the child to be baptized.
The babe died in his baptismal robes. Clovis blamed the baptism, but Clothilda rejoiced that God had
taken a soul directly to heaven to eternal bliss.

Another son was born, baptized, and fell sick. Clovis claimed that baptism would kill him, too, but
the mother prayed and he recovered. Then, in a battle with the Alamanni tribe, Clovis was facing a
terrible defeat. He cried, “Jesus Christ, Clothilda says thou art the Son of the living God, and
thou canst give victory to those who hope in thee. Give me victory and I will be baptized. I have
tried my gods and they have deserted me. I call on thee. Only save me.” The king of the Alamanni
fell, and his army fled. Clovis returned and told Clothilda.

She summoned the bishop of Rheims, who exhorted the king to renounce his gods. “Yes, holy father,”
said Clovis, “but my people will not consent. However, I will speak to them.” He did, and with one
accord they renounced their German gods. The baptistry was hung with tapestries, fragrant candles
flared, the aroma of incense filled the shrine with fragrance so that many thought they were amid
the odors of paradise. Clovis advanced like another Constantine to the baptismal font and the bishop
said, “Bend your neck. Worship what you burned and burn what you worshiped.” Thus the king was
baptized in the name of the Father and of the Son and of the Holy Spirit.

Three thousand of his army followed him in baptism.

This first of many mass conversions during the Middle Ages raises important questions about the
paganizing of Christianity. Missions can proceed in two ways. One is the way of individual
conversion with a period of instruction prior to baptism. In general this is the method used by
Protestant missions under the evangelical movements of the nineteenth century, with their emphasis
upon individual change of heart. The disadvantage of this method is that the Christian converts in a
pagan culture become, by reason of their change in faith, uprooted from their own culture and
compelled to move into an alien enclave.

The other method is mass conversion, and it was this method that converted Europe. Kings like Clovis
embraced the faith, and their people followed them into the spiritual City as into the earthly. This
meant that individuals were not uprooted from their culture, but it also meant that the converts
brought with them into the church their superstitions and behavior.

This is evident in many instances, starting with Clovis himself. Jesus was for him a tribal war-god.
The Franks especially admired St. Peter, whose noblest exploit in their eyes was his eagerness to
wield his sword to protect the Lord Jesus and to slice off the ear of the high priest’s servant.
This admiration for militant religion is also reflected in St. George, a military saint who became
the patron of England, and St. James, the patron saint of Christian Spain in the struggle against
Islam.

The saints, with their particular assignments, may well have meant more to the people than Christ
himself. St. Anthony took care of pigs, St. Gaul looked after hens, St. Apollonia, whose jaw had
been broken in the persecution, cured the toothache, St. Genevieve cured fever, and St. Blaise was
responsible for sore throats. For almost every human need these tenuously converted Germans could
find a saint.

Many tales circulated about the miraculous powers of the saints. The story was told of two beggars,
one lame, the other blind. They happened to be caught in a procession carrying the relics of St.
Martin and were fearful lest they be cured and so deprived of their alms. The one who could see but
not walk mounted the shoulders of the one who could walk but not see, and they hurried to get beyond
the range of the saint’s miraculous powers, but, poor fellows, they failed to make it.

\section{PASSAGE TO ENGLAND}

In spite of the shallow impact of the gospel among the Franks, their conversion did provide a
passage through Gaul to Great Britain. The invasion of Britain by the Anglo-Saxons (or English) made
hostilities so great that any idea of the Britons evangelizing the Germans was unthinkable. The
initiative for the conversion of the English came from Rome. In 596 Pope Gregory the Great (596–604)
sent a party of Benedictine monks to distant and barbaric England under the leadership of
Augustine—another St. Augustine, who would become known as St. Augustine of Canterbury.

Augustine and his monks began their ministry in Kent, one of the twelve areas controlled by the
Anglo-Saxon invaders of England. Under the favor of Queen Bertha, another of those Christian queens
eager to convert a pagan husband, Augustine secured a hearing from King Ethelbert, but only out of
doors where Augustine would be less able to exercise his magical powers. Rumor had it that he was
able to make tails grow on the backs of those who displeased him. The king was so persuaded by
Augustine that he granted land for the foundation of a monastery at Canterbury, ever after to be the
seat of the English religious leader. Gregory appointed Augustine the first archbishop of
Canterbury.

The missionaries who followed Augustine worked farther north. By the time of King Oswy in the late
seventh century, the two missionary thrusts converged, Celtic followers of Columba working toward
the south and those of Augustine toward the north. Oswy’s queen was from the south and followed the
Roman practices, but Oswy had received his Christian beliefs from the north and observed the Celtic
way. Among the points of dispute was the date for Easter. When the King had ended Lent and was
keeping Easter, the Queen and her household were still fasting: enough to ruin any family’s Easter!

At the Synod of Whitby in 664, Oswy brought the matter to a decision. The Celtic advocate appealed
to the authority of Columba, the Romans to that of Peter, to whom Christ gave the keys. “Is that
really so? Does Peter guard the gates to heaven?” Oswy asked the Celtic defender. He, of course,
agreed. Oswy promptly resolved to take no chances of alienating the doorkeeper of heaven! He agreed
to follow Roman practices. After Whitby the British Isles moved relentlessly into the orbit of Rome.

Once securely rooted in Anglo-Saxon England, Christianity returned in fresh vigor to the Continent.
The most famous British missionary of the Middle Ages was Winfrid, better known as Boniface. Born in
Devonshire, Boniface was commissioned by Pope Gregory II in 729 to evangelize Germany. His primary
task was to convert the pagan population, and in this he had great success. Besides that, however,
he brought the British and Irish missionary monks and their converts into closer relations with the
bishop of Rome.

Boniface was deeply devoted to the cause. He could not only chop down a sacred tree while the
shuddering multitude watched to see him struck dead by fire from heaven; he could also win the
masses to the faith and organize the districts in which he exercised his powers so as to leave
behind him a church structure bound firmly to the central authority at Rome.

Boniface became archbishop of Mainz and would have ended his splendid career there in peace if final
qualms about the failure of his early efforts in Frisia (the Netherlands) had not drawn him back to
that still pagan field. There he sealed his faith with martyrdom. If that was in the year 754, as is
generally supposed, the young Frankish prince Charles, who was destined to become the Emperor
Charlemagne and architect of a new—Christian—Europe, was twelve years old at the time.

\readinghead{Suggestions for Further Reading}
\begin{readinglist}
\entry Davis, R. H. C. A History of Medieval Europe. London: Longmans, 1957.
\entry Deanesley, Margaret. A History of the Medieval Church 590–1500. London: Methuen, 1969.
\entry Neill, Stephen. A History of Christian Missions. Middlesex: Penguin, 1964.
\entry Pirenne, Henri. A History of Europe from the Invasions to the XVI Century. New York: Murray
  Printing Company, 1938.
\entry *Southern, R. W. Church and Society in the Middle Ages. New Haven, CT: Yale University Press,
  1992.
\end{readinglist}
